\section{Reproduction and further uses}\label{nt:reproduction}
The accompanying runner executes both retained checkers with one command
from the repository root:
\begin{verbatim}
python3 papers/newton-calculus/check.py --output /tmp/newton-check
\end{verbatim}
The two checkers retain their distinct purposes.
The Newton checker compares physical matrix turns and cylindrical
traces with the displayed finite formulas in exact finite fields;
it includes the colored-circle coefficients and repeated-root
polynomial reduction. The support-and-profile checker enumerates
integer supports, checks different parenthesizations of multiple
products, and verifies rational instances of the stationary product
and harmonic identities. These finite tests detect coefficient,
indexing, phase and normalization errors. The all-level algebra,
interior maximum and limiting statements are proved in the text.

The positive-source limit uses a uniform factorial estimate, a
polynomial number of states and a positive summand approaching the
strict maximum at every sufficiently large level. Introducing retained
half twists or performing the fillings gives complex terms. The
resulting invariant requires separate bounds on all complementary
terms and a noncancellation argument. Likewise, stationary boundary
parameters are variational data; their global marked geometric
realization is not inferred here.

The polynomial relation in the finite calculation is the complete
one, $P_N$. Changing the main meridian changes that relation and some
turn factors. Those deformed identities require a separate calculation.

\begin{thebibliography}{9}
\bibitem{Finite}\emph{Finite spectral evaluation and uniform filling bounds}.
Companion paper, \href{https://sashakolpakov.github.io/murakamis_dream/papers/filling-bounds.pdf}{repository edition}.
\bibitem{Kashaev1995} R. M. Kashaev,
A link invariant from quantum dilogarithm,
\emph{Modern Physics Letters A} \textbf{10} (1995), 1409--1418.
\href{https://arxiv.org/abs/q-alg/9504020}{arXiv:q-alg/9504020}.
\bibitem{KauffmanLambropoulou} L. H. Kauffman and S. Lambropoulou,
On the classification of rational tangles,
\emph{Advances in Applied Mathematics} \textbf{33} (2004), 199--237.
\href{https://arxiv.org/abs/math/0311499}{arXiv:math/0311499}.
\bibitem{FuterGueritaud} D. Futer and F. Gu\'eritaud,
From angled triangulations to hyperbolic structures,
\emph{Contemporary Mathematics} \textbf{541} (2011), 159--182.
\href{https://arxiv.org/abs/1004.0440}{arXiv:1004.0440}.
\end{thebibliography}
\end{document}
